\section{Die Leiter}

\textbf{[B]} Die Leptonmassen lassen sich als Stufen einer $\xi$-Leiter
schreiben:
\[
  m_i = r_i\,\xi^{p_i}\,v ,
\]
mit rationalen Strukturzahlen $r_i$, rationalen Exponenten $p_i$ und einer
Brückenkonstante $v$. Für die Leptonen sind das
\[
  \bigl(r_e,p_e\bigr) = \Bigl(\tfrac43,\tfrac32\Bigr),\quad
  \bigl(r_\mu,p_\mu\bigr) = \Bigl(\tfrac{16}{5},1\Bigr),\quad
  \bigl(r_\tau,p_\tau\bigr) = \Bigl(\tfrac{25}{9},\tfrac23\Bigr).
\]

\textbf{[K]} In den \emph{Verhältnissen} fällt $v$ heraus, und es bleiben
parameterfreie Ausdrücke:
\[
  \frac{m_\mu}{m_e} = \frac{12}{5}\,\xi^{-1/2},\qquad
  \frac{m_\tau}{m_\mu} = \frac{125}{144}\,\xi^{-1/3},\qquad
  \frac{m_\tau}{m_e} = \frac{25}{12}\,\xi^{-5/6}.
\]

Nach A080 sind das die intrinsischen Größen: sie hängen an der Geometrie und an
nichts sonst. Genau an ihnen ist zu messen.

\section{Woher die Exponenten kommen}

\textbf{[K]} Die Exponenten $p_i$ sind nicht frei. Sie sind Windungszahlen auf
einem \emph{logarithmischen} Torus mit der Schrittweite $\Delta p = 1/3$. Diese
Schrittweite ist die fundamentale Terz der harmonischen Reihe -- dieselbe
Geometrie der periodischen Randbedingung, die in einer schwingenden Saite die
Obertöne erzeugt. Halbzahlige Werte sind geometrische Mitten zwischen zwei
Windungsebenen, negative Werte Gegenwindungen.

\textbf{[K]} Die drei Leptonexponenten liegen auf dieser Leiter:
\[
  p_e = \tfrac32,\qquad p_\mu = 1,\qquad p_\tau = \tfrac23,
\]
mit den Stufen
\begin{align}
  p_e   - p_\mu  &= \tfrac12 &&(\text{geometrische Mitte}), \notag\\
  p_\mu - p_\tau &= \tfrac13 &&(\text{fundamentale Terz}). \notag
\end{align}
Es sind dieselben Wicklungszahlen $(n_\theta, n_\phi)$, die auch die Vorfaktoren
$r_i = n_\phi^2/n_\theta$ festlegen (A110): Masse-Vorfaktor und
Exponent stammen aus einer Quelle, nicht aus zweien. Die Leiter ist damit keine
Anpassung, sondern die Auslesung der Resonanzstruktur des Torus.

\section{Was die Leiter liefert}

\textbf{[K]} Die Rechnung gegen die PDG-Werte ergibt:

\begin{center}
\begin{tabular}{lrrr}
\hline
Verhältnis & Leiter & PDG & Abweichung \\
\hline
$m_\mu/m_e$    & $207{,}85$ & $206{,}77$ & $+0{,}52\,\%$ \\
$m_\tau/m_\mu$ & $16{,}992$ & $16{,}817$ & $+1{,}04\,\%$ \\
$m_\tau/m_e$   & $3532$     & $3477$     & $+1{,}57\,\%$ \\
\hline
\end{tabular}
\end{center}

\textbf{[K]} Der Toleranzmaßstab nach A010 und A080 verlangt drei Angaben
nebeneinander. Die Messgenauigkeit der Vergleichsgrößen liegt bei $m_\mu/m_e$
im Bereich von $10^{-9}$, bei $m_\tau$ bei etwa $7\cdot10^{-5}$. Die
Bestimmungstoleranz der Leiterwerte ist um Größenordnungen gröber: die
Strukturzahlen sind rational deklariert, nicht auf Stellen bestimmt.

\textbf{Daraus folgt die Bewertung:} die Leiter trägt die \emph{Hierarchie} und
das Prozentniveau. Sie trägt keine Präzision. Eine Abweichung von einem Prozent
ist hier kein Befund gegen die Struktur, weil die Struktur auf diesem Niveau gar
nichts Genaueres behauptet -- aber sie ist auch kein Erfolg, den man auf fünf
Stellen ausweisen dürfte.

\section{Die Gegenrichtung: was die Daten für $\xi$ zurückgeben}

\textbf{[K]} Löst man die drei Verhältnisse nach $\xi$ auf, liefern die Daten
\[
  \xi_{\mu/e} = 1{,}347\cdot10^{-4},\quad
  \xi_{\tau/\mu} = 1{,}375\cdot10^{-4},\quad
  \xi_{\tau/e} = 1{,}358\cdot10^{-4},
\]
gegenüber $4/30000 = 1{,}3333\cdot10^{-4}$, also $+1{,}1$ bis $+3{,}2\,\%$.
Die drei Werte streuen untereinander um etwa $2\,\%$.

Das ist der ehrliche empirische Status: die Leiter ist in den Verhältnissen auf
dem Prozentniveau konsistent, nicht exakt. Sie gibt kein einheitliches $\xi$
zurück.

\textbf{[K]} Dasselbe von der anderen Seite: löst man je Lepton nach der
Brückenkonstante auf, ergibt sich $248{,}9$, $247{,}6$ und $245{,}1$ GeV,
gestreut um $246{,}22$ GeV. Die Restabweichung wird also in die
Brückenkonstante absorbiert und kürzt sich in den Verhältnissen \emph{nicht}
vollständig weg.

\section{Die Gegenprobe, die fehlschlägt}

\textbf{[K]} Die naheliegende Vermutung, die Restabweichungen folgten dem
fraktalen Korrekturfaktor, wird von der Rechnung \emph{verneint}.
$K_\text{frak} = 1-100\xi$ entspricht $-1{,}333\,\%$, die Halb- und
Dreihalbpotenzen $-0{,}665$ und $-1{,}995\,\%$. Die gemessenen Abweichungen
sind $+0{,}52$ und $+1{,}04\,\%$ -- sie passen weder im Betrag noch im
\emph{Vorzeichen} auf eine einzelne $K_\text{frak}$-Potenz.

Dieser Abschnitt steht hier, weil eine fehlgeschlagene Gegenprobe genauso zum
Befund gehört wie eine gelungene. Der naheliegende Erklärungsversuch ist
geprüft und gescheitert; er wird nicht durch einen freien Exponenten gerettet.

\textbf{Die auffällige Verdopplung ist arithmetisch erzwungen -- und belegt
keine Korrektur.} Die $\tau/\mu$-Abweichung ist fast genau das Doppelte der
$\mu/e$-Abweichung ($1{,}04\,\%$ gegen $0{,}52\,\%$), und die $\tau/e$-Abweichung
ist ihre Summe ($0{,}52 + 1{,}04 = 1{,}56\,\%$). Das ist kein neues Gesetz und
kein Zufall, sondern eine reine Buchhaltungsnotwendigkeit: entlang einer
multiplikativen Leiter, in der $m_\tau/m_e = (m_\tau/m_\mu)(m_\mu/m_e)$ gilt,
addieren sich die logarithmischen Abweichungen. Von den drei Abweichungen ist
daher nur \emph{eine} unabhängig gemessen (das Verhältnis $1{,}98$); die dritte
ist als Summe festgelegt.

\textbf{Wichtig ist, was dieser Schritt \emph{nicht} tut.} Die Additivität gilt
für \emph{jede} multiplikative Korrektur, unabhängig davon, welchen Wert oder
welches Vorzeichen sie hat. Sie ist deshalb kein Beleg für den fraktalen Faktor,
den die Gegenprobe gerade verworfen hat -- im Gegenteil, sie zeigt, dass die
Verdopplung überhaupt keine Korrektur auszeichnet: sie folgt schon aus der
Leiterform allein. Wer aus der Verdopplung ein Gesetz lesen will, liest eine
einzige gemessene Zahl dreifach.

\textbf{[K]} Der Faktor, der die Leiter pro Generation trägt (ihr
generationszählender Schritt, Größenordnung $N_0 \approx 39$), ist eine
\emph{andere} Größe als die feste Skala $100$ in $K_\text{frak}$: die eine ist
der Leiterschritt, die andere die RG-Skala der absoluten Korrektur (A040). Der
Wert dieses Leiterschritts ist nicht hergeleitet -- das ist die eine offene
Stelle, die weiter unten benannt wird.

\section{Zwei Rechenmodi, die sich ausschließen}

\textbf{[SETZUNG]} Bevor eine dieser Abweichungen gebucht wird, ist der
Rechenmodus festzulegen. Es gibt zwei, und sie schließen einander aus.

\textbf{Modus 1, referenzfrei.} Alle Verhältnisse werden rein aus $\xi$
gerechnet, ohne Anker. Dann stehen alle drei unabhängig gegen die Daten, jede
$0{,}5$ bis $1{,}6\,\%$ daneben, und alle drei sind legitime \emph{Prüfpunkte}.

Ein Wort zur Sprachregelung, weil hier sonst ein Widerspruch zu A010 entsteht:
diese drei Verhältnisse sind \emph{keine} Prognosegrößen der Serie. Die
Prognoseliste umfasst vier Größen, und die Leptonmassen erscheinen dort allein
über die $\tau$-Masse aus dem Zirkulanten (A110), nicht über die Leiter. Die
Leiter liefert Prüfpunkte auf Prozentniveau -- das ist weniger als eine
Prognose und mehr als nichts.

\textbf{Modus 2, verankert.} $m_\mu/m_e$ wird als Referenzpunkt deklariert.
Dann ist dieses Verhältnis \emph{verbraucht}: seine Abweichung darf nicht mehr
als Prüfergebnis gebucht werden, sonst wird dieselbe Größe zweimal benutzt --
einmal als Anker, einmal als Test.

\textbf{[B]} Die beiden Modi zu mischen ist die häufigste Fehlbuchung in diesem
Sektor. In der A-Serie wird durchgehend Modus 1 verwendet, sofern nichts anderes
ausdrücklich gesagt ist.

\section{Prüfskript}

\begin{itemize}
  \item Alle Zahlen dieses Dokuments -- Leiterverhältnisse, Reverse-$\xi$,
    Brückenkonstanten je Lepton, die fehlgeschlagene
    $K_\text{frak}$-Gegenprobe:
    \texttt{python/A\_Serie\_Skripte/a100\_leiter.py}
\end{itemize}


\section{Zeichenerklärung}

Symbole die in der gesamten A-Serie verwendet werden, sind in A010 erklärt.
Spezifisch für dieses Dokument:

\begin{center}
\renewcommand{\arraystretch}{1.3}
{\small\begin{tabular}{lp{9cm}}
\toprule
Symbol & Bedeutung \\
\midrule
$r_i\in\mathbb{Q}$ & Rationaler Vorfaktor in $m_i=r_i\cdot\xi^{p_i}\cdot v$ \\
$p_i\in\mathbb{Q}$ & Exponent; Leiter-Schrittweite $\Delta p=1/3$ \\
$(n,l,j)$ & Haupt-, Bahn-, Gesamtdrehimpuls der Mode \\
$v=246\,\text{GeV}$ & Higgs-Vakuumerwartungswert als Energieeinheit \\
\bottomrule
\end{tabular}}
\renewcommand{\arraystretch}{1.0}
\end{center}

\section{Was offen bleibt}

\textbf{Erstens ist die Prozentstreuung gewollt, kein Defekt.} Die Leiter ist
eine Aussage auf Prozentniveau; ein aus ihr zurückgerechnetes $\xi$ streut
entsprechend um $2\,\%$. Genau deshalb verankert der Sektor die absolute Skala
nicht in der Leiter, sondern in einem einzigen Referenzpunkt (A110): die Leiter
trägt die Hierarchie, der Referenzpunkt die Präzision. Die Streuung ist die
richtige Folge dieser Arbeitsteilung, nicht ein ungelöstes Problem.

\textbf{Zweitens sind die Exponenten $p_i$ hergeleitet, nicht bloß motiviert.}
Sie sind Windungszahlen auf dem logarithmischen Torus (oben), mit der
Schrittweite $1/3$ und geometrischen Mitten bei halbzahligen Werten -- dieselbe
Quelle, aus der die Vorfaktoren $r_i$ folgen. Damit ist die Bedingung aus A020
erfüllt: der Exponent liegt unabhängig fest, die $\xi$-Potenz ist eine Aussage.

\textbf{Drittens ist die Verdopplung geklärt.} Sie ist arithmetisch erzwungen
(oben): entlang einer multiplikativen Leiter addieren sich die logarithmischen
Abweichungen, nur eine ist unabhängig. Es bleibt kein ungebuchter Rest.

\textbf{Viertens bleibt genau ein Punkt offen -- und er ist benannt.} Die eine
unabhängig gemessene Abweichung entspricht einem generationszählenden
Leiterschritt der Größenordnung $N_0 \approx 39$, dessen Wert nicht hergeleitet
ist. Er ist von der festen RG-Skala $100$ in $K_\text{frak}$ zu trennen (A040)
und wird auch nicht durch die Brückenkonstante erklärt, in der die Restabweichung
absorbiert wird (A080). Der Wert dieses Schritts und die Kandidaten dafuer sind in A105 gesammelt; ob hinter ihm Struktur sitzt, ist die einzige wirklich offene Frage des Sektors.

\section{Ersetzt}

Dieses Dokument nimmt auf: Dok.~006 (Massenformel), Dok.~046 (Massenberechnung),
Dok.~016 (Teilchenmassen), Dok.~189 (Windungszahlen der Exponenten), Dok.~190
(R55 Additivität, P42 Referenzpunkt), sowie die Leiter-Abschnitte aus Dok.~292.
Eingearbeitet: K2 (Exponent $p_\tau$), P40 (Ebenentrennung), P42 (Referenzpunkt),
P35, R55 (Evidenzbreite: nur eine Stelle unabhängig gemessen).

\section{Verwendung von KI-Werkzeugen}

Dieses Dokument ist unter Verwendung KI-gestützter Werkzeuge entstanden. Die
Fragestellungen, die Setzungen und alle inhaltlichen Entscheidungen stammen vom
Autor; die Werkzeuge formulieren aus, rechnen nach und erstellen Prüfskripte.
Die Verantwortung für Inhalt und Fehler liegt vollständig beim Autor. Der
ausführliche Wortlaut steht in A010.
